\begin{proof}[Proof of \cref{obj:lem:published-cate-benchmark}]
\begin{enumerate}
\item The hypotheses give
\[
d>0,\qquad \gamma>0,\qquad \alpha+\beta>0.
\]
Define the real denominators
\[
D_{\mathrm{or}}:=2\gamma+d,
\qquad
D_{\mathrm{pair}}:=2\gamma+d+\frac{\gamma d}{\alpha+\beta}.
\]
Both are strictly positive:
\[
D_{\mathrm{or}}>0,\qquad D_{\mathrm{pair}}>0.
\]
% lean: CausalSmith.Stat.TwosamplePointcateAnnotationFrontier.published_cate_benchmark

\item Clearing the positive denominators and simplifying gives
\[
\frac{1}{1+d/(2\gamma)+d/(4s_{\mathrm{pub}})}
=\frac{2\gamma}{D_{\mathrm{pair}}},
\qquad
\frac{1}{2+d/\gamma}
=\frac{\gamma}{D_{\mathrm{or}}}.
\]
Likewise,
\[
\frac{d/4}{1+d/(2\gamma)}
=\frac{1}{2}\frac{\gamma d}{D_{\mathrm{or}}}.
\]
Since \(s_{\mathrm{pub}}=(\alpha+\beta)/2\), multiplication by \(2>0\) therefore yields
\[
s_{\mathrm{pub}}<\frac{d/4}{1+d/(2\gamma)}
\quad\Longleftrightarrow\quad
\alpha+\beta<\frac{\gamma d}{D_{\mathrm{or}}}.
\]
These are the asserted exponent identities and cutoff equivalence.
% lean: CausalSmith.Stat.TwosamplePointcateAnnotationFrontier.published_cate_benchmark

\item Fix \(n\in\mathbb N\) with \(n\ge2\), so that
\[
n\ge1.
\]
After substituting the exponent identities, the numerical piecewise benchmark becomes
\[
\begin{cases}
n^{-2\gamma/D_{\mathrm{pair}}},
& s_{\mathrm{pub}}<(d/4)/(1+d/(2\gamma)),\\
n^{-\gamma/D_{\mathrm{or}}},
& s_{\mathrm{pub}}\ge(d/4)/(1+d/(2\gamma)).
\end{cases}
\]
We compare these two powers by splitting at the equivalent cutoff in the preceding step.
% lean: CausalSmith.Stat.TwosamplePointcateAnnotationFrontier.published_cate_benchmark

\item Suppose first that
\[
\alpha+\beta<\frac{\gamma d}{D_{\mathrm{or}}}.
\]
Multiplying by \(D_{\mathrm{or}}>0\), then dividing by \(\alpha+\beta>0\), gives
\[
(\alpha+\beta)D_{\mathrm{or}}<\gamma d,
\qquad
D_{\mathrm{or}}<\frac{\gamma d}{\alpha+\beta}.
\]
Consequently,
\[
2D_{\mathrm{or}}
\le D_{\mathrm{or}}+\frac{\gamma d}{\alpha+\beta}
=D_{\mathrm{pair}}.
\]
Multiplying by \(\gamma>0\) and dividing by the positive product
\(D_{\mathrm{pair}}D_{\mathrm{or}}\) yields
\[
\frac{2\gamma}{D_{\mathrm{pair}}}
\le\frac{\gamma}{D_{\mathrm{or}}},
\qquad
-\frac{\gamma}{D_{\mathrm{or}}}
\le-\frac{2\gamma}{D_{\mathrm{pair}}}.
\]
For a base \(n\ge1\), real powers are nondecreasing in their exponent. Hence
\[
n^{-\gamma/D_{\mathrm{or}}}
\le n^{-2\gamma/D_{\mathrm{pair}}},
\]
and thus
\[
\max\left\{
n^{-\gamma/D_{\mathrm{or}}},
n^{-2\gamma/D_{\mathrm{pair}}}
\right\}
=n^{-2\gamma/D_{\mathrm{pair}}}.
\]
The cutoff equivalence selects precisely this first branch of the piecewise benchmark.
% lean: CausalSmith.Stat.TwosamplePointcateAnnotationFrontier.published_cate_benchmark

\item In the complementary case,
\[
\frac{\gamma d}{D_{\mathrm{or}}}\le\alpha+\beta.
\]
Multiplication by \(D_{\mathrm{or}}>0\), followed by division by
\(\alpha+\beta>0\), gives
\[
\gamma d\le(\alpha+\beta)D_{\mathrm{or}},
\qquad
\frac{\gamma d}{\alpha+\beta}\le D_{\mathrm{or}}.
\]
It follows that
\[
D_{\mathrm{pair}}
=D_{\mathrm{or}}+\frac{\gamma d}{\alpha+\beta}
\le2D_{\mathrm{or}}.
\]
Multiplying by \(\gamma>0\) and dividing by
\(D_{\mathrm{or}}D_{\mathrm{pair}}>0\) now gives
\[
\frac{\gamma}{D_{\mathrm{or}}}
\le\frac{2\gamma}{D_{\mathrm{pair}}},
\qquad
-\frac{2\gamma}{D_{\mathrm{pair}}}
\le-\frac{\gamma}{D_{\mathrm{or}}}.
\]
Monotonicity of real powers for \(n\ge1\) therefore implies
\[
n^{-2\gamma/D_{\mathrm{pair}}}
\le n^{-\gamma/D_{\mathrm{or}}},
\]
so
\[
\max\left\{
n^{-\gamma/D_{\mathrm{or}}},
n^{-2\gamma/D_{\mathrm{pair}}}
\right\}
=n^{-\gamma/D_{\mathrm{or}}}.
\]
The cutoff equivalence selects the second branch, including equality at the cutoff. Substituting the definitions of \(D_{\mathrm{or}}\) and \(D_{\mathrm{pair}}\) proves the asserted maximum identity in both cases.
% lean: CausalSmith.Stat.TwosamplePointcateAnnotationFrontier.published_cate_benchmark
\end{enumerate}
\end{proof}
